\documentclass{article}

\usepackage{arxiv}

\usepackage[spanish]{babel}
\usepackage[utf8]{inputenc}
\usepackage[T1]{fontenc}
\usepackage{amsmath}
\usepackage{amssymb}
\usepackage{hyperref}
\usepackage{cleveref}
\usepackage{url}
\usepackage{booktabs}
\usepackage{amsfonts}
\usepackage{nicefrac}
\usepackage{microtype}
\usepackage{graphicx}
\usepackage{float}
\usepackage{subcaption}
\crefname{table}{cuadro}{cuadros}
\Crefname{table}{Cuadro}{Cuadros}


\title{TP 2.2 - GENERADORES PSEUDOALEATORIOS DE DISTINTAS DISTRIBUCIONES DE PROBABILIDAD}


\author{
    Alvaro Tomasino \\
    Ingeniería en Sistemas de Información\\
    Universidad Tecnológica Nacional - FRRO\\
    Zeballos 1341, S2000, Argentina\\
    \texttt{atomasino08@gmail.com} \\
\And
    Juan Bautista Martinelli \\
    Ingeniería en Sistemas de Información\\
    Universidad Tecnológica Nacional - FRRO\\
    Zeballos 1341, S2000, Argentina\\
    \texttt{martinellij4@gmail.com} \\
\And
    Tomás Aguilera \\
    Ingeniería en Sistemas de Información\\
    Universidad Tecnológica Nacional - FRRO\\
    Zeballos 1341, S2000, Argentina\\
    \texttt{tomas.aguilera.27@gmail.com} \\
\And
    Melina Aronson \\
    Ingeniería en Sistemas de Información\\
    Universidad Tecnológica Nacional - FRRO\\
    Zeballos 1341, S2000, Argentina\\
    \texttt{melinaaronson@gmail.com} \\
}




\begin{document}

\maketitle




\begin{abstract}
El objetivo de este informe es, con los generadores usados en su anterior parte, empezar a
utilizarlos en distintos contextos para finalizar con nuestro estudio.
\end{abstract}

\keywords{Simulación \and generadores pseudoaleatorios \and distribuciones de probabilidad \and transformada inversa \and método de rechazo}




\section{Introducción}

Al emplear los generadores como los programamos anteriormente, estos van a generar
valores uniformes continuos entre 0 y 1.

Entonces ¿Cómo hacemos para generar valores de distintas distribuciones? La solución es
apoyarse en un material elaborado por Thomas Naylor para redescubrir el mecanismo de
los generadores, y poder implementarlo en los casos que se desea y de distintas
probabilidades.

A continuación, durante todo el informe, se desarrollan las bases teóricas de las siguientes
distribuciones, con el desarrollo de los diferentes métodos de generación, y su
implementación y test en código Python según indica la siguiente tabla.

\begin{table}[H]
    \centering
    \begin{tabular}{lccccc}
        \toprule
        \textbf{Distribución de Probabilidad} & \textbf{Tipo} & \textbf{T. Inversa} & \textbf{M. Rechazo} & \textbf{Código} & \textbf{Testeo} \\
        \midrule
        UNIFORME          & continua & sí & sí & sí & sí \\
        EXPONENCIAL       & continua & sí & sí & no & sí \\
        GAMMA             & continua & sí & no & no & no \\
        NORMAL            & continua & sí & sí & sí & sí \\
        PASCAL            & discreta & sí & no & no & no \\
        BINOMIAL          & discreta & sí & no & no & sí \\
        HIPERGEOMÉTRICA   & discreta & sí & no & no & no \\
        POISSON           & discreta & sí & no & no & sí \\
        EMPÍRICA DISCRETA & discreta & sí & no & no & sí \\
        \bottomrule
    \end{tabular}
    \caption{Listado de distribuciones y tareas a realizar con cada una.}
    \label{tab:distribuciones}
\end{table}




\section{Metodología}
\label{sec:metodologia}

\subsection{Distribución de Probabilidades (deducción matemática)}

\subsubsection{Distribución de Probabilidad Uniforme}

Es una familia de distribuciones de probabilidad para variables aleatorias, tales que para
cada miembro de la familia todos los intervalos de igual longitud en la distribución en su
rango son igualmente probables. El dominio está definido por dos parámetros,
\textbf{a} y \textbf{b}, que son sus valores mínimo y máximo respectivamente.

\paragraph{Notación}

Si $X$ es una variable aleatoria continua con distribución uniforme continua se nota como:

\begin{equation}
    X \sim U(a, b) \quad \text{o} \quad X \sim \mathrm{Uniforme}(a, b)
\end{equation}

\paragraph{Función de densidad}

Si $X \sim U(a, b)$ entonces la función de densidad es:

\begin{equation}
    f_X(x) = \frac{1}{b - a} \quad \forall\, x \in [a, b]
\end{equation}


\subsubsection{Distribución de Probabilidad Exponencial}

La distribución exponencial es una distribución continua que se utiliza para modelar
tiempos de espera para la ocurrencia de un cierto evento, siendo un caso particular de la
distribución gamma.
Esta distribución al igual que la distribución geométrica tiene la propiedad de pérdida de
memoria.

\paragraph{Notación}

Si $X$ es una variable aleatoria continua con distribución exponencial continua, y $\lambda > 0$
se nota como:

\begin{equation}
    X \sim \mathrm{Exp}(\lambda) \quad \text{con } \lambda > 0
\end{equation}

\paragraph{Función de densidad}

Si $X \sim \mathrm{Exp}(\lambda)$ entonces la función de densidad es:

\begin{equation}
    f_X(x) = \lambda\, e^{-\lambda x} \quad \forall\, x \geq 0
\end{equation}


\subsubsection{Distribución de Probabilidad Gamma}

La distribución Gamma es una distribución continua utilizada para modelar tiempos de
espera, confiabilidad de sistemas y diversos fenómenos donde se acumulan eventos
aleatorios.

Su distribución depende de dos parámetros:

\begin{itemize}
    \item $\boldsymbol{\alpha}$ \textbf{(alfa)}: parámetro de forma (\textit{shape}), $\alpha > 0$.
    \item $\boldsymbol{\beta}$ \textbf{(beta)}: parámetro de escala (\textit{scale}), $\beta > 0$.
\end{itemize}

Siendo su denotación:

\begin{equation}
    X \sim \Gamma(\alpha, \beta)
\end{equation}

\paragraph{Función de Densidad de Probabilidad}

\begin{equation}
    f(x) = \frac{1}{\Gamma(\alpha)\,\beta^{\alpha}}\, x^{\alpha - 1}\, e^{-x/\beta}
\end{equation}

Para $x > 0$, donde:

\begin{equation}
    \Gamma(\alpha) = \int_0^{\infty} t^{\alpha - 1} e^{-t}\, dt
\end{equation}

\paragraph{Obtención de la Función Acumulada}

La función acumulada se define como:

\begin{equation}
    F(x) = P(X \leq x)
\end{equation}

Por definición: $F(x) = \int f(t)\,dt$.

Sustituyendo la Función de Densidad:

\begin{equation}
    F(x) = \int_0^{x} \frac{1}{\Gamma(\alpha)\,\beta^{\alpha}}\, t^{\alpha - 1}\, e^{-t/\beta}\, dt
\end{equation}

Sacando las constantes:

\begin{equation}
    F(x) = \frac{1}{\Gamma(\alpha)\,\beta^{\alpha}} \int_0^{x} t^{\alpha - 1} e^{-t/\beta}\, dt
\end{equation}

Realizando el cambio de variables:

\begin{equation}
    u = \frac{t}{\beta}, \quad t = \beta u, \quad dt = \beta\, du
\end{equation}

Queda:

\begin{equation}
    F(x) = \frac{1}{\Gamma(\alpha)} \int_0^{x/\beta} u^{\alpha - 1} e^{-u}\, du
\end{equation}

Llamando a la integral $\gamma(\alpha, z)$, quedaría finalmente:

\begin{equation}
    F(x) = \frac{\gamma\!\left(\alpha,\, x/\beta\right)}{\Gamma(\alpha)}
\end{equation}


\subsubsection{Distribución de Probabilidad Normal}

La distribución normal es una distribución continua que aparece con más frecuencia en
estadística. La gráfica de su función de densidad tiene una forma acampanada y es
simétrica respecto a un determinado parámetro estadístico. Esta curva se conoce como
campana de Gauss.

\paragraph{Notación}

Si $X$ es una variable aleatoria continua con distribución normal, se nota como:

\begin{equation}
    X \sim N(\mu,\, \sigma^2)
\end{equation}

donde los parámetros de entrada son la media ($\mu$) y la varianza ($\sigma^2$).

\paragraph{Función de densidad}

La función de densidad de probabilidad está dada por la ecuación de la campana de Gauss:

\begin{equation}
    f(x) = \frac{1}{\sigma\sqrt{2\pi}}\, e^{-\frac{1}{2}\left(\frac{x - \mu}{\sigma}\right)^2}
\end{equation}

Para el caso de la distribución normal estándar ($\mu = 0$ y $\sigma = 1$), la fórmula se simplifica a:

\begin{equation}
    f(z) = \frac{1}{\sqrt{2\pi}}\, e^{-z^2/2}
\end{equation}

\paragraph{Obtención de la función acumulada}

La función de distribución acumulada calcula la probabilidad de que una variable aleatoria
continua $X$ sea menor o igual a un valor específico $x$:

\begin{equation}
    F(x) = P(X \leq x)
\end{equation}

Para una distribución normal, se define a través de la integral de su función de densidad de
probabilidad:

\begin{equation}
    F(x) = \frac{1}{\sigma\sqrt{2\pi}} \int_{-\infty}^{x} e^{-\frac{1}{2}\left(\frac{t - \mu}{\sigma}\right)^2} dt
\end{equation}


\subsubsection{Distribución de Probabilidad Pascal}

También conocida como \textbf{distribución binomial negativa}, es una distribución discreta que
modela el número de fracasos que ocurren antes de obtener un número fijo de éxitos en
una secuencia de ensayos de Bernoulli independientes.

\paragraph{Parámetros}

La distribución depende de dos parámetros:

\begin{itemize}
    \item $r$: número fijo de éxitos que se desean obtener ($r > 0$).
    \item $p$: probabilidad de éxito en cada ensayo ($0 < p < 1$).
\end{itemize}

Denotada como:

\begin{equation}
    X \sim \mathrm{Pascal}(r, p)
\end{equation}

donde $X$ representa el número de fracasos antes de alcanzar el $r$-ésimo éxito.

\paragraph{Función de Probabilidad}

La probabilidad de obtener exactamente $x$ fracasos antes del $r$-ésimo éxito es:

\begin{equation}
    P(X = x) = \binom{x + r - 1}{r - 1} p^r (1-p)^x \quad \text{para } x = 0, 1, 2, \ldots
\end{equation}

\paragraph{Obtención de la Función Acumulada}

Se define como: $F(x) = P(X \leq x)$. Por definición:

\begin{equation}
    F(x) = \sum_{k=0}^{x} \binom{k + r - 1}{r - 1} p^r (1-p)^k
\end{equation}

Y desarrollando algunos términos:

\begin{equation}
    F(x) = p^r + r\,p^r(1-p) + \binom{r+1}{r-1} p^r(1-p)^2 + \cdots
\end{equation}

En forma general:

\begin{equation}
    F(x) = \sum_{k=0}^{x} \binom{k + r - 1}{r - 1} p^r (1-p)^k
\end{equation}

\paragraph{Obtención de la Función Inversa}

Partiendo de $u = F(x)$ con $u \sim U(0, 1)$, debemos encontrar el menor valor entero $x$
tal que:

\begin{equation}
    F(x) \geq u \qquad \text{es decir:} \quad x = F^{-1}(u)
\end{equation}


\subsubsection{Distribución de Probabilidad Binomial}

La distribución binomial es una distribución de probabilidad discreta que cuenta el número
de éxitos en una secuencia de $n$ ensayos de Bernoulli independientes entre sí con una
probabilidad fija $p$ de ocurrencia de éxito entre los ensayos. Un experimento de Bernoulli
se caracteriza por ser dicotómico, es decir, solo dos resultados son posibles. A uno se lo
llama ``éxito'' y tiene una probabilidad de ocurrencia $p$, y al otro se lo llama ``fracaso'' y
tiene una probabilidad $q = 1 - p$.

\paragraph{Notación}

Si una variable aleatoria discreta $X$ tiene una distribución binomial con parámetros
$n \in \mathbb{N}$ y $p$ con $0 < p < 1$, se nota como:

\begin{equation}
    X \sim \mathrm{Bin}(n, p)
\end{equation}

\paragraph{Función de densidad}

Su función de densidad está dada por la función de probabilidad:

\begin{equation}
    P[X = x] = \binom{n}{x} p^x (1-p)^{n-x} \quad \text{para } x = 0, 1, 2, \ldots, n
\end{equation}

\paragraph{Obtención de la función acumulada}

La función de distribución acumulada está dada por:

\begin{equation}
    F_X(x) = P[X \leq x] = \sum_{k=0}^{\lfloor x \rfloor} \binom{n}{k} p^k (1-p)^{n-k}
    \quad \text{si } 0 \leq x \leq n
\end{equation}


\subsubsection{Distribución de Probabilidad Hipergeométrica}

Es una distribución discreta que modela la probabilidad de obtener una determinada
cantidad de éxitos al extraer elementos de una población finita sin reemplazo.

\paragraph{Parámetros}

La distribución depende de 3 parámetros:

\begin{itemize}
    \item $N$: tamaño total de la población.
    \item $K$: número de elementos considerados éxitos dentro de la población.
    \item $n$: tamaño de la muestra extraída.
\end{itemize}

Se denota:

\begin{equation}
    X \sim H(N, K, n)
\end{equation}

donde $X$ es el número de éxitos obtenidos en la muestra.

\paragraph{Función de Probabilidad}

Se define como:

\begin{equation}
    P(X = x) = \frac{\dbinom{K}{x}\dbinom{N-K}{n-x}}{\dbinom{N}{n}}
\end{equation}

donde $\binom{K}{x}$ representa las formas de elegir $x$ éxitos, $\binom{N-K}{n-x}$ las formas de
elegir los fracasos y $\binom{N}{n}$ el total de muestras posibles, con:

\begin{equation}
    \max(0,\, n - (N - K)) \leq x \leq \min(n, K)
\end{equation}

\paragraph{Obtención de la Función Acumulada}

Se define como: $F(x) = P(X \leq x)$. Por definición:

\begin{equation}
    F(x) = \sum_{k=0}^{x} P(X = k)
\end{equation}

Sustituyendo:

\begin{equation}
    F(x) = \sum_{k=0}^{x} \frac{\dbinom{K}{k}\dbinom{N-K}{n-k}}{\dbinom{N}{n}}
\end{equation}

Finalmente:

\begin{equation}
    F(x) = \frac{1}{\dbinom{N}{n}} \sum_{k=0}^{x} \binom{K}{k}\binom{N-K}{n-k}
\end{equation}


\subsubsection{Distribución de Probabilidad Poisson}

La distribución de Poisson es una distribución de probabilidad discreta que expresa, a
partir de una frecuencia de ocurrencia media, la probabilidad de que ocurra un determinado
número de eventos durante cierto periodo de tiempo. Se especializa en la probabilidad de
ocurrencia de sucesos con probabilidades muy pequeñas, o sucesos raros.

\paragraph{Notación}

Sea $\lambda > 0$ y $X$ una variable aleatoria discreta; si $X$ tiene una distribución de Poisson
con parámetro $\lambda$, se nota como:

\begin{equation}
    X \sim \mathrm{Poisson}(\lambda) \quad \text{o} \quad X \sim \mathrm{Poi}(\lambda)
\end{equation}

\paragraph{Función de densidad}

Como la variable aleatoria $X$ es discreta, no usa una función de densidad, sino una función
de masa de probabilidad:

\begin{equation}
    P[X = k] = \frac{e^{-\lambda}\,\lambda^k}{k!}
\end{equation}

donde $k$ es el número de ocurrencias del evento, y $\lambda$ representa el número de veces
que se espera que ocurra dicho fenómeno durante un intervalo dado.

\paragraph{Obtención de la función acumulada}

La función acumulada calcula la probabilidad de que ocurran $x$ eventos o menos en un
intervalo dado. Se obtiene sumando las probabilidades puntuales desde 0 hasta el valor
de $x$:

\begin{equation}
    P(X \leq x) = \sum_{k=0}^{x} \frac{e^{-\lambda}\,\lambda^k}{k!}
\end{equation}


\subsubsection{Distribución de Probabilidad Empírica Discreta}

En algunas situaciones, ninguna de las distribuciones teóricas conocidas describe bien el
comportamiento de los datos observados. En esos casos se usan las distribuciones
empíricas, que permiten aproximar el comportamiento de un fenómeno a partir de
frecuencias de datos recolectados directamente de la realidad.

\paragraph{Notación}

Sea $X$ una variable aleatoria discreta que puede tomar un conjunto de valores específicos
$b_1, b_2, \ldots, b_n$. Asociada a cada valor, existe una probabilidad observada $p_i$ tal que:

\begin{equation}
    P(X = b_i) = p_i
\end{equation}

\paragraph{Función de probabilidad}

La distribución se define por una tabla de frecuencias donde se listan los posibles
resultados y sus probabilidades de ocurrencia. La suma de todas las probabilidades debe
ser igual a la unidad:

\begin{equation}
    \sum_{i=1}^{n} p_i = 1
\end{equation}


\subsection{Método de la Transformada Inversa (deducción matemática)}

El Método de la Transformada Inversa es una técnica para generar números aleatorios que
sigan una distribución de probabilidad determinada a partir de números aleatorios uniformes
en el intervalo $(0,1)$. Es decir, tenemos:

\begin{equation}
    r \sim U(0, 1)
\end{equation}

Donde $r$ es un número entre 0 y 1 generado a partir de una distribución Uniforme.

Queremos obtener una variable aleatoria $X$ con una distribución específica arbitraria, cuya
función de distribución acumulada es $F(x)$. La función acumulada $F(x)$ indica la
probabilidad de que la variable $X$ tome un valor menor o igual que $x$:

\begin{equation}
    F(x) = \int_{-\infty}^{x} f(t)\, dt \qquad F(x) = P(X \leq x)
\end{equation}

Esta función de probabilidad acumulada está definida en el intervalo de 0 a 1.

El método consiste en aplicar la función inversa de la distribución acumulada para que $x$
tenga un valor dado por la distribución deseada. Esto es, si generamos un valor uniforme
$r$ entre 0 y 1 y buscamos el valor $x$ tal que:

\begin{equation}
    r = F(x) = \int_{-\infty}^{x} f(t)\, dt
\end{equation}

entonces:

\begin{equation}
    F^{-1}(r) = x
\end{equation}

$F^{-1}(r)$ es la transformación inversa o mapeo de $r$ sobre el intervalo o dominio de $x$.

Gráficamente, dado el valor $r_0$ en el eje vertical de la función acumulada $F(x) = r$, su
correspondiente $x_0$ en el eje horizontal es el valor generado con la distribución deseada.
El criterio $P(X \leq x) = F(x) = P[r \leq F(x)] = P[F^{-1}(r) \leq x]$ equivale a resolver la
ecuación en $x$ en términos de $r$.

Este método cuenta con varias desventajas que se van a ir presentando a medida que lo
vayamos aplicando a las diversas distribuciones. La principal desventaja es que no siempre
existe una inversa analítica de la función de distribución acumulada $F(x)$, por ejemplo la
distribución Gamma o la distribución normal estándar. Esto se resuelve aplicando métodos
numéricos. Cuando la inversa no existe en forma explícita, hay que calcularla mediante
algoritmos iterativos. Esto tiene un costo computacional elevado. Evaluar funciones
logarítmicas, exponenciales o resolver ecuaciones numéricamente puede ser más lento que
otros métodos de generación.

A continuación se deduce el método matemático para generar números aleatorios a partir de
una distribución Uniforme para cada distribución mencionada en la sección anterior.


\subsubsection{Distribución de Probabilidad Uniforme}

Dada la función de densidad $f(x)$ de la distribución Uniforme $X \sim U(0, 1)$:

\begin{equation}
    f(x) = \frac{1}{b - a}, \quad a \leq x \leq b
\end{equation}

con función de distribución acumulada $F(x)$:

\begin{equation}
    F(x) = \int_a^x \frac{1}{b-a}\, dt = \frac{x - a}{b - a}
\end{equation}

Dado un número aleatorio $r$ generado por $r \sim U(0, 1)$, igualamos $F(x) = r$:

\begin{equation}
    \frac{x - a}{b - a} = r
\end{equation}

Buscamos $F^{-1}(r) = x$:

\begin{equation}
    x = a + (b - a)\,r
\end{equation}

Si $r \sim U(0,1)$, entonces $F^{-1}(r) = a + (b-a)r = x$ tiene exactamente la función de
distribución $F$. Para la distribución $U(0,1)$, este caso es especialmente limpio porque
$F(x)$ es lineal y su inversa se obtiene con álgebra directa, sin necesidad de métodos
numéricos.


\subsubsection{Distribución Exponencial}

Dada la función de densidad $f(x)$ de la distribución Exponencial:

\begin{equation}
    f(x) = \lambda\, e^{-\lambda x}, \quad x \geq 0
\end{equation}

con función de distribución acumulada $F(x)$:

\begin{equation}
    F(x) = \int_0^x \lambda\, e^{-\lambda t}\, dt = \left[-e^{-\lambda t}\right]_0^x = 1 - e^{-\lambda x}
\end{equation}

Despejando $x$ a partir de $F(x) = r$:

\begin{align}
    1 - e^{-\lambda x} &= r \notag \\
    e^{-\lambda x}     &= 1 - r \notag \\
    -\lambda x         &= \ln(1 - r) \notag \\
    x                  &= -\frac{1}{\lambda}\ln(1 - r)
\end{align}

Como $r \sim U(0, 1)$, entonces $(1 - r)$ también es $U(0, 1)$. Por lo tanto se puede escribir
directamente:

\begin{equation}
    x = -\frac{1}{\lambda}\ln(r)
\end{equation}


\subsubsection{Distribución Gamma}

Dada la función de densidad $f(x)$ de la distribución Gamma:

\begin{equation}
    f(x) = \frac{1}{\beta^k\,\Gamma(k)}\, x^{k-1}\, e^{-x/\beta}, \quad x \geq 0
\end{equation}

donde:

\begin{equation}
    \Gamma(k) = \int_0^{\infty} t^{k-1} e^{-t}\, dt
\end{equation}

con función de distribución acumulada $F(x)$:

\begin{equation}
    F(x) = \int_0^x \frac{1}{\beta^k\,\Gamma(k)}\, t^{k-1}\, e^{-t/\beta}\, dt
\end{equation}

Debido a que para una distribución Gamma no se puede formular explícitamente una
función de distribución acumulada, debemos considerar un método alternativo para generar
valores de variable aleatoria con esta distribución.

Igualar $F(x) = r$ y despejar $x$ \textbf{no tiene solución algebraica} en el caso general.
No existe una fórmula cerrada del tipo $x = \text{fórmula}(r, k, \beta)$. Esto significa que
$F^{-1}(r)$ debe obtenerse \textbf{numéricamente}, lo cual hace que la transformada inversa
directa sea costosa computacionalmente.

Cuando $k$ no es entero, las opciones son:

\begin{itemize}
    \item \textbf{Métodos numéricos} para invertir $F(x)$ iterativamente (bisección, Newton-Raphson).
    \item \textbf{Algoritmos especializados} como el de Cheng (1977), basado en el método de
    rechazo, que es el más usado en la práctica para $k > 1$.
\end{itemize}

Cuando $k$ es un entero positivo, existe un camino alternativo muy elegante. Se puede
demostrar que si:

\begin{equation}
    X_1, X_2, \ldots, X_k \sim \mathrm{Exponencial}(\lambda) \quad \text{independientes}
\end{equation}

entonces su suma:

\begin{equation}
    X = X_1 + X_2 + \cdots + X_k \sim \mathrm{Gamma}\!\left(k,\, \beta = \tfrac{1}{\lambda}\right)
\end{equation}

A esto se le conoce como \textbf{distribución de Erlang}. Aplicando la transformada inversa de
cada exponencial (que sí tiene forma cerrada):

\begin{equation}
    X = -\frac{1}{\lambda}\sum_{i=1}^{k} \ln(r_i) = -\frac{1}{\lambda}\ln\!\left(\prod_{i=1}^{k} r_i\right)
\end{equation}


\subsubsection{Distribución Normal}

Dada la función de densidad $f(x)$ de la distribución Normal:

\begin{equation}
    f(x) = \frac{1}{\sigma\sqrt{2\pi}}\, e^{-\frac{1}{2}\left(\frac{x-\mu}{\sigma}\right)^2},
    \quad -\infty < x < \infty
\end{equation}

con función de distribución acumulada $F(x)$:

\begin{equation}
    F(x) = \int_{-\infty}^{x} \frac{1}{\sigma\sqrt{2\pi}}\, e^{-\frac{1}{2}\left(\frac{t-\mu}{\sigma}\right)^2} dt
\end{equation}

Al igual que la Gamma, esta integral no tiene forma cerrada. Se expresa en términos de la
función error:

\begin{equation}
    F(x) = \frac{1}{2}\left[1 + \mathrm{erf}\!\left(\frac{x - \mu}{\sigma\sqrt{2}}\right)\right]
\end{equation}

Por lo tanto $F^{-1}(r)$ tampoco tiene forma algebraica directa, y el camino de la
transformada inversa pura requeriría métodos numéricos.

Las soluciones a esto son el método de rechazo para la Normal, que opera sobre la normal
estándar y usa como función envolvente una distribución que es fácil de muestrear
(típicamente una combinación de uniformes y exponenciales), y el \textbf{Método de la suma de
uniformes} (Teorema Central del Límite).

El Teorema Central del Límite establece que la suma de muchas variables aleatorias
independientes con cualquier distribución tiende a una normal. Entonces, si sumamos $K$
uniformes$(0,1)$:

\begin{equation}
    S = r_1 + r_2 + \cdots + r_K
\end{equation}

para $K$ suficientemente grande, $S$ se comporta aproximadamente como una normal.

Cada $r_i \sim \mathrm{Uniforme}(0,1)$ tiene:

\begin{equation}
    \theta = \frac{a + b}{2} = \frac{1}{2}, \qquad \sigma = \frac{b - a}{\sqrt{12}} = \frac{1}{\sqrt{12}}
\end{equation}

Por el TLC, la suma de $K$ de ellas tiene:

\begin{equation}
    E\!\left(\sum r_i\right) = \frac{K}{2}, \qquad \mathrm{Var}\!\left(\sum r_i\right) = \frac{K}{12}
\end{equation}

Estandarizando (restando la media y dividiendo por el desvío):

\begin{equation}
    z = \frac{\displaystyle\sum_{i=1}^{K} r_i - K/2}{\sqrt{K/12}}
\end{equation}

Este $z$ es aproximadamente $N(0,1)$, y luego $x = \mu + \sigma z$ da cualquier normal deseada.

Por definición $z$ es un valor de variable aleatoria con distribución normal estándar que se
puede escribir de la forma:

\begin{equation}
    z = \frac{x - \mu_x}{\sigma_x}
\end{equation}

Igualando las últimas dos ecuaciones obtenemos:

\begin{equation}
    \frac{x - \mu_x}{\sigma_x} = \frac{\displaystyle\sum_{i=1}^{K} r_i - K/2}{\sqrt{K/12}}
\end{equation}

y resolviendo para $x$, se tiene que:

\begin{equation}
    x = \sigma_x \left(\frac{12}{K}\right)^{1/2} \!\left(\sum_{i=1}^{K} r_i - \frac{K}{2}\right) + \mu_x
\end{equation}

Así se pueden generar valores aleatorios con distribución normal con promedio $\mu_x$ y
desvío $\sigma_x$ a partir de $K$ valores generados por una distribución uniforme.


\subsubsection{Distribución Pascal}

Dada la función de densidad $f(x)$ de la distribución Pascal:

\begin{equation}
    P(X = x) = \binom{x + r - 1}{r - 1} p^r (1-p)^x
\end{equation}

con función de distribución acumulada $F(x)$:

\begin{equation}
    F(x) = P(X \leq x) = \sum_{k=s}^{x} \binom{k - 1}{s - 1} p^s (1-p)^{k-s}
\end{equation}

Como $F(x)$ es discreta, no se puede despejar $x$ algebraicamente. En su lugar, dado
$r \sim \mathrm{Uniforme}(0,1)$, se busca el menor $x$ tal que $F(x) \geq r$:

\begin{equation}
    F(x-1) < r \leq F(x)
\end{equation}

Es decir, se recorre la función acumulada de izquierda a derecha hasta que supera a $r$ por
primera vez. Su algoritmo es el siguiente:

\begin{enumerate}
    \item Comenzar en $x = s$ (el mínimo posible).
    \item Calcular $P(X = x)$ con la fórmula.
    \item Acumular $F(x) = F(x-1) + P(X = x)$.
    \item Si $F(x) \geq r$, devolver $x$ como valor generado.
    \item Si no, incrementar $x = x + 1$ y volver al paso 2.
\end{enumerate}


\subsubsection{Distribución Binomial}

Dada la función de densidad $f(x)$ de la distribución Binomial:

\begin{equation}
    P[X = x] = \binom{n}{x} p^x (1-p)^{n-x}
\end{equation}

con función de distribución acumulada $F(x)$:

\begin{equation}
    F(x) = \sum_{k=0}^{x} \binom{n}{k} p^k (1-p)^{n-k}
\end{equation}

Tomamos $r \sim U(0,1)$. Como $F(x)$ es discreta, la función inversa $F^{-1}(x)$ se obtiene
por iteración (búsqueda). Se aplica la misma regla que en la Pascal: se define como el
menor valor entero $x$ que cumple la condición $F(x) \geq r$:

\begin{equation}
    F(x-1) < r \leq F(x)
\end{equation}

El algoritmo es el siguiente:

\begin{enumerate}
    \item Generar $r \sim \mathrm{Uniforme}(0,1)$.
    \item Iniciar en $x = 0$ con $F(x) = 0$.
    \item En cada paso sumar $P(X = x)$ a $F(x)$.
    \item Si $F(x) \geq r$, devolver $x$.
    \item Si no, incrementar $x$ y repetir desde el paso 3.
\end{enumerate}


\subsubsection{Distribución Hipergeométrica}

Dada la función de densidad $f(x)$ de la distribución Hipergeométrica:

\begin{equation}
    P(X = x) = \frac{\dbinom{K}{x}\dbinom{N-K}{n-x}}{\dbinom{N}{n}},
    \quad x = \max(0, n + K - N), \ldots, \min(n, K)
\end{equation}

con función de distribución acumulada $F(x)$:

\begin{equation}
    F(x) = \sum_{k=x_{\min}}^{x} \frac{\dbinom{K}{k}\dbinom{N-K}{n-k}}{\dbinom{N}{n}}
\end{equation}

Tomamos $r \sim \mathrm{Uniforme}(0,1)$. Como $F(x)$ es discreta, se aplica la misma regla
que en Binomial y Pascal:

\begin{equation}
    F(x-1) < r \leq F(x)
\end{equation}

La búsqueda termina en $x = \min(n, K)$ como máximo, ya que $F(\min(n,K)) = 1 \geq r$
siempre. El algoritmo es el siguiente:

\begin{enumerate}
    \item Generar $r \sim \mathrm{Uniforme}(0,1)$.
    \item Calcular $x_{\min} = \max(0, n + K - N)$ e iniciar con $F(x) = 0$.
    \item En cada paso sumar $P(X = x)$ a $F(x)$.
    \item Si $F(x) \geq r$, devolver $x$.
    \item Si no, incrementar $x$ y repetir desde el paso 3.
\end{enumerate}


\subsubsection{Distribución de Poisson}

Dada la función de densidad $f(x)$ de la distribución de Poisson:

\begin{equation}
    P(X = x) = \frac{e^{-\lambda}\,\lambda^x}{x!}, \quad x = 0, 1, 2, \ldots
\end{equation}

con función de distribución acumulada $F(x)$:

\begin{equation}
    F(x) = \sum_{k=0}^{x} \frac{e^{-\lambda}\,\lambda^k}{k!}
\end{equation}

Tomamos $r \sim \mathrm{Uniforme}(0,1)$. Como $F(x)$ es discreta, se aplica la misma regla
que en las distribuciones discretas anteriores. El algoritmo es el siguiente:

\begin{enumerate}
    \item Generar $r \sim \mathrm{Uniforme}(0,1)$.
    \item Iniciar en $x = 0$ con $F(x) = 0$.
    \item En cada paso sumar $P(X = x)$ a $F(x)$.
    \item Si $F(x) \geq r$, devolver $x$.
    \item Si no, incrementar $x$ y repetir desde el paso 3.
\end{enumerate}


\subsubsection{Distribución Empírica Discreta}

Dada la función de densidad $f(x)$ de la distribución Empírica Discreta:

\begin{equation}
    P(X = x_i) = p_i, \quad i = 1, 2, \ldots, n
\end{equation}

con la condición de que:

\begin{equation}
    \sum_{i=1}^{n} p_i = 1, \quad p_i \geq 0
\end{equation}

con función de distribución acumulada $F(x)$:

\begin{equation}
    F(x_i) = \sum_{k=1}^{i} p_k
\end{equation}

Tomamos $r \sim \mathrm{Uniforme}(0,1)$. Se aplica la misma regla que en todas las
discretas anteriores. El algoritmo es el siguiente:

\begin{enumerate}
    \item Generar $r \sim \mathrm{Uniforme}(0,1)$.
    \item Iniciar en $i = 1$ con $F = 0$.
    \item En cada paso sumar $p_i$ a $F$.
    \item Si $F \geq r$, devolver $x_i$.
    \item Si no, incrementar $i$ y repetir desde el paso 3.
\end{enumerate}


\subsection{Método de rechazo (deducción matemática)}

El objetivo es generar una variable aleatoria $X$ que siga una función de densidad de
probabilidad objetivo $f(x)$.

Si $f(x)$ es una función acotada y $x$ tiene además un rango finito con $a \leq x \leq b$,
entonces se puede utilizar la técnica de rechazos para generar los valores de variables
aleatorios. La aplicación de esta técnica requiere que proceda de acuerdo con las
siguientes etapas:

\begin{enumerate}
    \item Normalizar el rango de $f$ mediante un factor de escala $c$ tal que
          $c \cdot f(x) \leq 1$, $a \leq x \leq b$.
    \item Definir a $x$ como una función lineal de $r$:
          \begin{equation}
              x = a + (b - a)\,r
          \end{equation}
    \item Generar parejas de números aleatorios $(r_1, r_2)$.
    \item Siempre que se encuentre una pareja de números aleatorios que satisfagan la
          relación
          \begin{equation}
              r_2 \leq c \cdot f\!\left[a + (b-a)\,r_1\right]
          \end{equation}
          dicho par será aceptado y se utilizará $x = a + (b-a)\cdot r_1$ como el valor
          generado de la variable aleatoria.
\end{enumerate}

La teoría sobre la que se apoya este método se basa en el hecho ya conocido relativo a la
probabilidad de que $r$ sea menor o igual a $c \cdot f(x)$:

\begin{equation}
    P\!\left[r \leq c \cdot f(x)\right] = c \cdot f(x)
\end{equation}

En consecuencia, si $x$ se elige al azar dentro del rango $(a, b)$ de acuerdo con la ecuación
$x = a + (b-a)r$ y en el caso de que $r > c \cdot f(x)$ se rechaza, la función de densidad
de probabilidad de los valores de $x$ aceptados deberá ser igual a $f(x)$.

Tocher ha demostrado que la esperanza matemática del número de intentos que se
realizan, antes de encontrar una pareja exitosa es igual a $1/c$. Esto implica que para
ciertas funciones de densidad de probabilidad este método puede resultar sumamente
ineficaz.


\subsubsection{Distribución de Probabilidad Uniforme}

Para aplicar el método de rechazo a la distribución uniforme continua, partimos de su
función de densidad de probabilidad objetivo, la cual está definida en el intervalo cerrado
$[a, b]$:

\begin{equation}
    f(x) = \frac{1}{b - a} \quad \forall\, x \in [a, b]
\end{equation}

\begin{enumerate}
    \item Buscamos un factor de escala $c$ tal que el máximo de la curva escalada sea
          menor o igual a 1: $c \cdot f(x) \leq 1$.

          Sustituyendo la función de densidad: $c \cdot \dfrac{1}{b-a} \leq 1$.

          Para optimizar el algoritmo y hacer que el ``techo'' sea exactamente 1,
          despejamos el factor de escala óptimo:
          \begin{equation}
              c = b - a
          \end{equation}

          Si comprobamos la función escalada con este valor de $c$, obtenemos:
          \begin{equation}
              c \cdot f(x) = (b - a) \cdot \frac{1}{b - a} = 1
          \end{equation}

    \item Mapeamos el primer número pseudoaleatorio $r_1$ al rango deseado:
          \begin{equation}
              x = a + (b - a) \cdot r_1
          \end{equation}

    \item La pareja de números pseudoaleatorios $(r_1, r_2)$ será aceptada si se cumple
          la relación $r_2 \leq c \cdot f(x)$. Reemplazando por los valores obtenidos en el
          primer paso, la condición se reduce a:
          \begin{equation}
              r_2 \leq 1
          \end{equation}
\end{enumerate}

Dado que por definición cualquier número pseudoaleatorio $r_2$ generado por la computadora
pertenece al intervalo $[0, 1)$, la inecuación se cumplirá en el 100\,\% de los casos.

Esta deducción demuestra matemáticamente que aplicar el método de rechazo para
generar valores de una distribución uniforme resulta en una redundancia donde no se
rechaza a ningún candidato. La eficiencia del algoritmo es máxima, comprobando
empíricamente el teorema de Tocher.


\subsubsection{Distribución Exponencial}

Para aplicar el método de rechazo a la distribución exponencial continua, partimos de su
función de densidad de probabilidad objetivo. Dado que la técnica de rechazos requiere que
la variable tenga un rango finito $[a, b]$, definimos obligatoriamente el límite inferior $a = 0$
(ya que no existen ocurrencias negativas) y truncamos la cola infinita en un límite superior
$b$ lo suficientemente grande para contener la mayor parte de la probabilidad:

\begin{equation}
    f(x) = \lambda\, e^{-\lambda x} \quad \forall\, x \in [0, b]
\end{equation}

\begin{enumerate}
    \item Buscamos un factor de escala $c$ tal que el máximo de la curva escalada sea
          menor o igual a 1: $c \cdot f(x) \leq 1$.

          Como la función exponencial decrece monotónamente, su valor máximo absoluto
          se encuentra siempre en el límite inferior, es decir, cuando $x = 0$:
          \begin{equation}
              f(0) = \lambda\, e^{-\lambda \cdot 0} = \lambda
          \end{equation}

          Para optimizar el algoritmo y hacer que el ``techo'' superior sea exactamente 1,
          determinamos que el factor de escala óptimo es:
          \begin{equation}
              c = \frac{1}{\lambda}
          \end{equation}

          Si comprobamos la función escalada con este valor de $c$, obtenemos la
          simplificación matemática clave de este método:
          \begin{equation}
              c \cdot f(x) = \frac{1}{\lambda} \cdot \lambda\, e^{-\lambda x} = e^{-\lambda x}
          \end{equation}

    \item Mapeamos el primer número pseudoaleatorio $r_1$ al rango finito que
          establecimos. Sabiendo que $a = 0$, la ecuación general se reduce a:
          \begin{equation}
              x = b \cdot r_1
          \end{equation}

    \item Siempre que se encuentre una pareja de números pseudoaleatorios $(r_1, r_2)$
          que satisfaga la relación $r_2 \leq c \cdot f(x)$, dicho par será aceptado.
          Reemplazando por la función normalizada obtenida en el primer paso, la condición
          final queda definida como:
          \begin{equation}
              r_2 \leq e^{-\lambda x}
          \end{equation}

          O expresado en función exclusivamente de la pareja de pseudoaleatorios
          generada:
          \begin{equation}
              r_2 \leq e^{-\lambda(b \cdot r_1)}
          \end{equation}
\end{enumerate}

A diferencia de lo demostrado con la distribución uniforme, aquí la inecuación no se
cumplirá en todos los casos. A medida que el candidato generado $r_1$ acerque la
coordenada $x$ hacia el límite $b$, el valor del techo $e^{-\lambda x}$ decaerá rápidamente
asintóticamente hacia cero, aumentando de forma exponencial la probabilidad de que el
número pseudoaleatorio $r_2$ sea mayor y, por ende, rechazado.

Esto ilustra el postulado de Tocher, demostrando que la ineficacia del método de rechazo
aumenta drásticamente cuando el área de la caja delimitadora es significativamente mayor
que el área bajo la curva de la función de densidad acotada.


\subsubsection{Distribución de Probabilidad Normal}

Para aplicar el método de rechazo a la distribución normal, partimos de su función de
densidad de probabilidad objetivo. Dado que esta distribución se define en un rango infinito
$(-\infty, \infty)$, para su implementación mediante esta técnica hay que truncar la
distribución en un intervalo finito $[a, b]$ lo suficientemente amplio para asegurar que la
masa de probabilidad fuera de ese rango sea despreciable:

\begin{equation}
    f(z) = \frac{1}{\sqrt{2\pi}}\, e^{-z^2/2}
\end{equation}

\begin{enumerate}
    \item Buscamos un factor de escala $c$ tal que el máximo de la curva escalada sea
          menor o igual a 1: $c \cdot f(x) \leq 1$.

          Como el valor máximo de la función normal estándar ocurre en la media ($z = 0$)
          y es aproximadamente $0{,}3989$, el factor de escala óptimo para que el ``techo''
          de la función sea exactamente 1 es:
          \begin{equation}
              c = \sqrt{2\pi} \approx 2{,}5066
          \end{equation}

    \item Se genera un primer número pseudoaleatorio $r_1 \sim U(0,1)$ y se proyecta
          linealmente al rango de interés $[a, b]$ para obtener un valor candidato $z$:
          \begin{equation}
              z = a + (b - a)\,r_1
          \end{equation}

    \item Se genera un segundo número pseudoaleatorio $r_2 \sim U(0,1)$. La pareja de
          números $(r_1, r_2)$ es aceptada si se satisface la relación de la función objetivo
          escalada:
          \begin{equation}
              r_2 \leq c \cdot f(z)
          \end{equation}

          Sustituyendo los valores de la distribución normal, la condición de aceptación se
          simplifica:
          \begin{equation}
              r_2 \leq e^{-z^2/2}
          \end{equation}

    \item Si la inecuación se cumple, el par es aceptado y se usa $z$ como el valor generado
          de la variable aleatoria. Si no, el candidato se rechaza y se repite el procedimiento
          desde el paso 2 hasta obtener un valor válido.
\end{enumerate}


\subsection{Código Método Rechazo}

\subsubsection{Distribución de Probabilidad Uniforme}

% Insertar código Python del método de rechazo para la distribución uniforme.

\subsubsection{Distribución de Probabilidad Normal}

% Insertar código Python del método de rechazo para la distribución normal.




\section{Gráficas}
\label{sec:graficas}

Se usaron diagramas de dispersión para ver el comportamiento de los generadores
basados en el método de rechazo (Uniforme, Normal).

% Insertar gráficas correspondientes.




\section{Discusión}
\label{sec:discusion}

% Insertar discusión de los resultados obtenidos.




\section{Conclusiones}
\label{sec:conclusiones}

% Insertar conclusiones del trabajo.




\nocite{*}

\bibliographystyle{unsrturl}
\bibliography{references}

\end{document}